\section{A plane inside a four-dimensional module}

Joining two two-dimensional modules requires a coefficient operation
that is absent when one constituent has dimension one. Both constituents
carry a nonzero holomorphic potential. Their critical coefficients
must be multiplied before the result can be pushed through the space
of extensions. The relevant operation compares the external product
of two vanishing-cycle complexes with the vanishing cycles of the sum
of their functions.

Take the three-loop matrix potential
\[
 V_n=\Mat_n(\C)^3,\qquad f_n(A,B,C)=\Tr(A[B,C]),
 \qquad Z_n=\Crit(f_n),\qquad 1\le n\le4.
\]
The commutator is $[B,C]=BC-CB$. The nondegenerate trace pairing
identifies the three derivatives with the three commutators, so
\[
 Z_n=\{[A,B]=[B,C]=[C,A]=0\}.
\]
A commuting triple is a representation of $\C[x,y,z]$; a one-dimensional
representation is its evaluation at a point of $\C^3$.
Let $E_{2,2}$ be the space of triples in $V_4$ preserving a plane
$H\subset\C^4$, together with that plane. In a basis adapted to $H$,
\[
 A=\begin{pmatrix}A_1&u_A\\0&A_2\end{pmatrix},\quad
 B=\begin{pmatrix}B_1&u_B\\0&B_2\end{pmatrix},\quad
 C=\begin{pmatrix}C_1&u_C\\0&C_2\end{pmatrix}.
\]
Matrix multiplication gives the decisive identity
\[
                  f_4|_{E_{2,2}}=f_2(A_1,B_1,C_1)+f_2(A_2,B_2,C_2).
\]
Neither term is zero as a function. For example, with
$D=\operatorname{diag}(1,-1)$, $U=e_{12}$ and $L=e_{21}$,
$\Tr(D[U,L])=2$. Taking this triple in each block gives $f_4=4$.
Taking its negative in the first matrix of the second block gives
$f_4=0$ while the two block potentials are $2$ and $-2$.
Thus the zero fibre of the sum is larger than the product of zero
fibres. The coefficient comparison must specify its restriction.

All sheaves and cochains below have coefficients in $\Q$. A
constructible complex has finite-dimensional stalk cohomology locally
constant on each piece of a finite algebraic stratification. We work
in the derived category, identifying maps up to quasi-isomorphism;
a quasi-isomorphism induces an isomorphism on every stalk cohomology
group. Our shift satisfies $K[r]^j=K^{j+r}$ and has differential
$(-1)^r d_K$. The tensor differential and interchange are
\[
 d(x\otimes y)=dx\otimes y+(-1)^{|x|}x\otimes dy,
 \qquad x\otimes y\longmapsto(-1)^{|x||y|}y\otimes x.
\]
Only ordinary cohomology and algebraic tensor products are used.

Changes of basis are retained by
$\frX_n=[V_n/\GL_n]$ and $\frM_n=[Z_n/\GL_n]$.
A constructible complex on this quotient means compatible complexes
on the whole action diagram
$V_n,\GL_n\times V_n,\GL_n^2\times V_n,\ldots$, with its composition
identifications. In particular a scalar triple retains its
$\GL_n$ stabilizer. All atlas pullbacks are ordinary pullbacks,
without a relative-dimension shift.

\section{Vanishing cycles as cochains with support}
\label{ts22:sec:support}

For $g:X\to\C$, put $X_0=g^{-1}(0)$ and write $i:X_0\hookrightarrow X$.
Nearby cycles $\psi_gK$ take cochains on the inverse image of the
universal cover of a small punctured disk in the value of $g$, then
restrict toward $X_0$. Restriction gives $i^*K\to\psi_gK$.
We use normalized vanishing cycles
\[
 \Phi_gK=\Fib(i^*K\to\psi_gK)=\Cone(i^*K\to\psi_gK)[-1].
\]
The fibre of $u:A\to B$ has terms $A^j\oplus B^{j-1}$ and
$d(a,b)=(da,u(a)-db)$. This fixes the sign as well as the shift.
If $K$ is supported on $X_0$, then $\Phi_gK=i^*K$.
In particular $\Phi_0=1$.
For a smooth ambient space and a small local Milnor fibre $F_{g,x}$,
the fibre sequence gives
\[
 H^j(\Phi_g\Q)_x=\widetilde H^{j-1}(F_{g,x};\Q).
\]
Here $F_{g,x}=B_\delta(x)\cap g^{-1}(\epsilon)$ with
$0<|\epsilon|\ll\delta\ll1$ in a range of constant local topology;
reduced cohomology removes its constant class.

For a closed subset $P\subset X$, let $R\Gamma_PK$ denote the
complex of sections with support in $P$. With $j:X\setminus P\hookrightarrow X$,
it is the fibre of $K\to Rj_*j^*K$. Choose the positive real ray
in the value of $g$ and set $P_g=\{\Re g\le0\}$. Then
\[
             \Phi_gK\simeq i^*R\Gamma_{P_g}K.
\]
To explain this identification, use a small tube over a disk in the
value of $g$. The part over $\Re g>0$ fibres over a contractible
half-disk and has the cochains of one nearby fibre. Its restriction
from the tube is the map used to define the fibre complex. Shrinking the tube
and the ambient neighborhood gives the stalk identification and its
restriction maps. A continuous rotation of the ray yields the
nearby-cycle continuation; a full positive rotation gives potential
monodromy. Product charts give the same construction after smooth
pullback. For the determinant and quadratic functions used below,
the required tubes and boundary-preserving transports are supplied
explicitly in Section~\ref{ts22:sec:wall}.

The normalized coefficient on the critical stack is
\[
 \cP_n=(\Phi_{f_n}\Q_{\frX_n}[2n^2])|_{\frM_n}.
\]
The shift is the ambient stack dimension $3n^2-n^2$.
Thus $\cP_1=\Q_{\C^3\times B\C^*}[2]$, while the ordinary pullback
of $\cP_2$ to $Z_2$ has shift eight.
The function $f_n$ has only critical value zero: its critical equations
force every commutator to vanish. Away from its critical locus it is
one smooth coordinate, so its vanishing cycles are zero.

\begin{lemma}[External cochains with support]
\label{ts22:lem:support-product}
For closed subsets $P\subset X$, $Q\subset Y$ locally defined by
finitely many real-analytic inequalities, and constructible complexes
$A,B$, there is a natural product isomorphism
\[
 R\Gamma_PA\boxtimes R\Gamma_QB
       \simeq R\Gamma_{P\times Q}(A\boxtimes B).
\]
It is associative under the tensor associator and is compatible with
restriction, smooth pullback, and proper direct image in either factor.
\end{lemma}

\begin{proof}
On small product neighborhoods, sections with support are relative
cochains for $(U,U\setminus P)$ and $(W,W\setminus Q)$. Their product
is the relative complex of
\[
 (U\times W,(U\setminus P)\times W\ \cup\ U\times(W\setminus Q)).
\]
The latter complement is exactly that of $P\times Q$.
Triangulate the pairs and their coefficient strata locally. The
oriented subdivision of a product simplex into shuffles gives the
relative cochain product; subdivision and the prism homotopy give
its inverse up to quasi-isomorphism. Over $\Q$, tensoring introduces
no torsion term. The boundary of a product cell is
$\partial(\sigma\times\tau)=\partial\sigma\times\tau+
(-1)^{\dim\sigma}\sigma\times\partial\tau$, which gives precisely
the stated tensor signs. Triple product cells give the two products
with their canonical associator. Refining triangulations preserves
the maps in the derived category, so these local descriptions patch.

Restriction and smooth product pullback preserve the pairs and their
product cells. For a proper map, cochains in inverse images of
shrinking neighborhoods compute cochains near the compact fibre:
a finite subcover shows that these neighborhoods are cofinal.
The same argument applies to the complementary open sets and their
relative complexes. It proves both proper base change and the
asserted proper compatibility of the product maps.
\end{proof}

There is therefore a canonical morphism on $X_0\times Y_0$,
\[
 \TS_{f,g}:\Phi_fA\boxtimes\Phi_gB
       \longrightarrow k^*\Phi_{f+g}(A\boxtimes B),
 \qquad k:X_0\times Y_0\hookrightarrow(f+g)^{-1}(0).
\]
It is the product of support complexes followed by the inclusion
$P_f\times P_g\subset P_{f+g}$. This defines the map independently
of any proof that it is an isomorphism. Rotating both half-planes
shows that it intertwines $T_f\boxtimes T_g$ with $T_{f+g}$.
If one potential is identically zero, its half-plane is the entire
factor and this is the ordinary smooth product identification.

\section{The comparison for two rank-two trace potentials}
\label{ts22:sec:wall}

The local geometry required by $\TS_{f,g}$ can be derived from
three traceless matrices. Write $A=aI+A_0$, and similarly for $B,C$.
In the basis
\[
 H=\begin{pmatrix}1&0\\0&-1\end{pmatrix},\quad E=e_{12},\quad F=e_{21}
\]
of $\mathfrak{sl}_2$, let $M$ have columns the coordinates of
$A_0,B_0,C_0$. Direct expansion gives
\[
                      f_2=2\det M.
\]
Indeed the trilinear form is alternating and its value on $(H,E,F)$
is two. Its critical equations are the nine cofactors, so its
critical set is $\operatorname{rank}M\le1$, with three free scalar
coordinates. At a nonzero rank-one point choose a nonzero pivot.
Block elimination gives
\[
 \det\begin{pmatrix}d&r\\c&D\end{pmatrix}
       =d\det(D-cd^{-1}r).
\]
The four entries of $D-cd^{-1}r$ are normal coordinates. Rescaling
one row absorbs the invertible factor $2d$. The resulting exact
normal form is a nondegenerate four-variable quadratic. At the
zero matrix the full cubic determinant remains. At rank at least
two the determinant is a submersion.
For a nondegenerate quadratic $q=\sum_{i=1}^m z_i^2$, its local
fibre over $\epsilon>0$ retracts onto $S^{m-1}$. Write $z=x+iy$:
then $x\cdot y=0$ and $|x|^2-|y|^2=\epsilon$. Shrink $y$ toward
zero and replace $x$ by
$\sqrt{\epsilon+|y|^2}\,x/|x|$. The norm decreases throughout,
so this retraction stays inside the local ball. This proves the
quadratic cohomology used below.

\begin{lemma}[The half-plane comparison]
\label{ts22:lem:wall}
For two copies of $f_2$ with constant rational coefficients,
$\TS_{f_2,f_2}$ is an isomorphism on $Z_2\times Z_2$. The same is
true after smooth pullback and after adjoining free coordinates.
\end{lemma}

\begin{proof}
Fix a point of $Z_2\times Z_2$ and write $u=\Re f_2$, $v=\Re f_2'$
for the two real values. The two support regions in the value plane
are the quadrant $C=\{u\le0,v\le0\}$ and the half-plane
$D=\{u+v\le0\}$. The comparison is induced by $C\subset D$.
Its being a quasi-isomorphism is equivalent, through the two
relative-cochain fibre sequences, to the inclusion
\[
   \{u+v>0\}\hookrightarrow\{u>0\text{ or }v>0\}
\]
inducing an equivalence on cochains in shrinking neighborhoods.

Here is an explicit deformation of the value regions. In the mixed
quadrant $u>0,v<0$ increase $v$ linearly to $\max(v,-u/2)$; in
$u<0,v>0$ increase $u$ linearly to $\max(u,-v/2)$.
On the positive quadrant and the positive axes use the identity.
These formulas agree continuously at the axes. They send the union
to $u+v>0$, never cross either zero axis in a mixed quadrant, and
preserve the subregion $u+v>0$ throughout the deformation. They
shrink the absolute value of the coordinate they change. Keep the
imaginary parts of both functions fixed.

This value deformation lifts in a cofinal system of small product
tubes. For completeness, the needed tube transversality can be
checked in each of the local forms just derived. A smooth function
is one coordinate. For the quadratic $\sum_{i=1}^4z_i^2$, a critical
point of the norm on a nonzero fibre satisfies
$z_i$ equal to a common phase times real numbers; hence
$\sum|z_i|^2=|\sum z_i^2|$. Therefore the sphere of radius $\delta$
is transverse over $0<|g|<\delta^2$. For $\det M$ on invertible
matrices, singular-value decomposition shows that a critical point
of the squared norm on a fixed determinant fibre must have all
three singular values equal. To derive this, vary two singular
values by $(\sigma_i e^t,\sigma_j e^{-t})$; the derivative of the
norm vanishes only if $\sigma_i=\sigma_j$. Conversely this is the
only possible norm-critical value, namely $3|\det M|^{2/3}$.
A sphere of radius $\delta$ is consequently transverse whenever
$0<|\det M|<(\delta^2/3)^{3/2}$.

Choose such balls, smaller value disks, and their products with
balls in free coordinates. Over a nonzero value the derivative of
the function is surjective both in the interior and on the boundary
of these balls. Local right inverses of that derivative lift a
value vector field; a partition of unity patches them to a field
tangent to the boundary. On each compact tube over a compact set of
nonzero values, its flow exists for the required time and preserves
the boundary. The preceding deformation stays in one mixed quadrant
until its endpoint and changes only the factor with a negative real
value. Thus it uses these lifts only at nonzero function values.
The positive coordinate remains fixed. For each point in a mixed
quadrant, the changed real coordinate stays in a compact interval
away from zero. The local horizontal transports therefore exist
through the whole deformation and depend continuously on that point.
At a positive axis the deformation is the identity in a neighborhood:
if $u>0$ is fixed near $u_0>0$ and $|v|<u/2$, the formula does not
move $v$, and the analogous statement holds on the other axis.
Consequently the lifted homotopies extend there by the identity.
The origin of the value plane is outside both complements, so no
transport through a singular zero fibre is required. The boundary-
preserving lifts stay inside the chosen product balls, while the
value paths stay in their disks because their moduli decrease.


The lifted homotopies identify the cochains of the two complementary
regions in this cofinal system. Hence their relative complexes, and
therefore the stalks of $\TS_{f,g}$, are quasi-isomorphic.
The map being tested was the canonical inclusion of supports, so
this stalk calculation proves that map is a sheaf isomorphism; no
choice of lifted homotopy enters its definition. Free-coordinate
balls and smooth pullback leave the argument unchanged.
\end{proof}

\begin{proposition}[Naturality needed for flag refinement]
\label{ts22:prop:coherence}
The support comparison commutes with proper pushforward and oriented
Gysin maps in either factor. Its two composites for three functions,
restricted to the product of their zero fibres, agree under the
ordinary tensor associator. These assertions hold as maps, without
requiring either composite to be an isomorphism.
\end{proposition}

\begin{proof}
For three functions, both composites start with support
$P_f\times P_g\times P_h$ and end with support $P_{f+g+h}$.
One factors the inclusion through $P_{f+g}\times P_h$, the other
through $P_f\times P_{g+h}$. The maps of sections with support for
nested closed sets compose as their inclusions do. The associative
product of relative cochains in Lemma~\ref{ts22:lem:support-product}
then proves the equality. No permutation of inputs occurs.

For a proper map $q:Y\to X$ with potential $g\circ q$, the inverse
image of $P_g$ is $P_{g\circ q}$. The proper comparison of relative
cochains in Lemma~\ref{ts22:lem:support-product} therefore identifies
$\Phi_gRq_*K$ with $Rq_{0*}\Phi_{g\circ q}K$, and identifies the
corresponding support-inclusion maps. For an embedding of complex
codimension $c$, its Thom map is cup product with the oriented
normal class of degree $2c$, followed by extension from the normal
neighborhood. This commutes with restriction of supports and with
the external product, with sign $(-1)^{2c|x|}=1$.
For integration over a compact complex flag fibre, evaluation on
its fundamental cycle commutes with the same restrictions and
products. These are the two operations forming every Gysin map
below. They establish the asserted compatibility before taking
vanishing cycles, and the fibre construction preserves it.
\end{proof}

\section{The coefficient for a two-by-two block split}
\label{ts22:sec:product}

For $d+e=n$, the incidence $E_{d,e}$ is a vector bundle of rank
$3(n^2-de)$ over the Grassmannian of $d$-planes, of dimension $de$.
It is smooth of dimension $3n^2-2de$. Its map $q_{d,e}$ to $V_n$
is proper because the flag variety is projective. The embedding
into $V_n\times\operatorname{Gr}(d,n)$ is the transverse zero set
of the three lower block entries, a normal bundle of rank $3de$.
Its positive complex Thom class, followed by flag integration,
defines
\[
 \gamma_{d,e}:Rq_{d,e*}\Q_{E_{d,e}}\longrightarrow\Q_{V_n}[4de].
\]
Normal integration gives degree $6de$ and the flag fibre subtracts
$2de$. For $(d,e)=(2,2)$ the dimensions are $40\to48$, and the
Gysin degree is sixteen.

Pass to $\cE_{d,e}=[E_{d,e}/\GL_n]$. Recording the two diagonal
blocks gives a smooth map
$a^{\mathrm{amb}}_{d,e}:\cE_{d,e}\to\frX_d\times\frX_e$.
In frames for the two blocks, the upper entries form an affine
space of dimension $3de$. Changes of splitting form the additive
group $\operatorname{Hom}(\C^e,\C^d)$ of dimension $de$.
Thus the relative stack dimension is $2de$. The function is pulled
back from $f_d+f_e$, by block multiplication.

The critical correspondence is the full fibre product
\[
 \cF_{d,e}=\cE_{d,e}\times_{\frX_n}\frM_n,
 \qquad
 \frM_d\times\frM_e\xleftarrow{\ a_{d,e}\ }\cF_{d,e}
                  \xrightarrow{\ b_{d,e}\ }\frM_n.
\]
Submodules and quotients of a commuting triple commute. This gives
$a_{d,e}$; it is not assumed smooth. Requiring only commuting
diagonal blocks would omit equations on the off-diagonal entries,
and would define a larger source.

Applying $\TS_{f,g}$ and smooth ambient pullback, with shifts,
gives on $\cF_{2,2}$
\[
 a_{2,2}^*(\cP_2\boxtimes\cP_2)
 \xrightarrow{\ \sim\ }
 \bigl(\Phi_{f_4|\cE_{2,2}}\Q_{\cE_{2,2}}[16]\bigr)|_{\cF_{2,2}}.
\]
The restriction to the product of zero fibres is legitimate here:
both diagonal blocks of a point of $\cF_{2,2}$ are critical.
It is precisely the restriction specified by $k$ in $\TS_{f,g}$.
Apply proper direct image, its vanishing-cycle exchange, and
$\Phi_{f_4}(\gamma_{2,2}[16])$. Then restrict to $\frM_4$.
This defines the actual coefficient morphism
\[
 \mu_{2,2}:Rb_{2,2*}a_{2,2}^*(\cP_2\boxtimes\cP_2)
                         \longrightarrow\cP_4.
\]
Its shifts are $16+16=32$; it has degree zero. The same construction
with one zero rank-one potential defines $\mu_{2,1}$, $\mu_{3,1}$,
and $\mu_{1,1}$, with Gysin degrees eight, twelve, and four.

These maps exist on the quotient stacks. At every action-diagram
level, all normal bundles pull back with their complex orientations;
face and degeneracy maps preserve the potentials and flag cycles.
The product of relative cochains and the inclusion of supports are
natural on these whole diagrams. Their product-cell and tubular
homotopies agree on iterated overlaps. This constructs descent of
the maps and their composition identifications, including stabilizers.
The ordinary cohomological multiplication $m_{d,e}$ is external
product, pullback by $a_{d,e}$, and application of $\mu_{d,e}$.

\section{The first associativity comparison with a nonzero pair}
\label{ts22:sec:assoc}

A flag of type $(2,1,1)$ is $H\subset K\subset\C^4$ with dimensions
two and three. Let $E_{211}$ be triples preserving this flag.
The flag variety has dimension $2\cdot1+2\cdot1+1\cdot1=5$.
Each preserving matrix has eleven entries, so
\[
 \dim E_{211}=38,\qquad f_4|_{E_{211}}=f_2|_H,
 \qquad \deg\gamma_{211}=2(48-38)=20.
\]
Its critical fibre product $\cF_{211}$ has maps $b_{211}$ to
$\frM_4$ and $a_{211}$ to $\frM_2\times\frM_1\times\frM_1$.
The constant-zero-factor comparison identifies its incoming
coefficient with the pullback of $\Phi_{f_2}\Q[12]$ on the ambient
flag space: the rank-two shift eight and the two rank-one shifts
two sum to twelve. Applying the degree-twenty Gysin map after
vanishing cycles gives
\[
 M_{211}:Rb_{211*}a_{211}^*(\cP_2\boxtimes\cP_1\boxtimes\cP_1)
                  \longrightarrow\cP_4.
\]
Unlike a complete flag of one-dimensional quotients, this flag
space has a nonzero potential.

Forgetting $H$ factors the ambient map through $E_{3,1}$.
The refinement has degree eight and the outer map degree twelve.
Forgetting $K$ factors it through $E_{2,2}$, where the refinement
has degree four and the outer map degree sixteen. Their shifts are
\[
            12\longrightarrow20\longrightarrow32,
       \qquad 12\longrightarrow16\longrightarrow32.
\]
The first refinement is the smooth ambient pullback of the
$(2,1)$ incidence on the three-dimensional subspace; the second is
the smooth ambient pullback of the $(1,1)$ incidence on the
two-dimensional quotient. Their critical Cartesian squares are
\[
\begin{array}{ccc}
 \cF_{211}&\longrightarrow&\cF_{2,1}\times\frM_1\\
 \downarrow&&\downarrow b_{2,1}\times1\\
 \cF_{3,1}&\longrightarrow&\frM_3\times\frM_1
\end{array}
 \qquad
\begin{array}{ccc}
 \cF_{211}&\longrightarrow&\frM_2\times\cF_{1,1}\\
 \downarrow&&\downarrow1\times b_{1,1}\\
 \cF_{2,2}&\longrightarrow&\frM_2\times\frM_2.
\end{array}
\]
In the first square the extra datum is a plane inside the invariant
three-space; in the second it is a line in the invariant rank-two
quotient. These descriptions prove that the squares retain the
same modules, flags, and changes of frame.

Proper base change and the product of cochains identify the common
source of $M_{211}$ with the source of each parenthesized expression.
Denote these identifications by $C_L,C_R$. They give the maps
\[
\begin{split}
 M_L&=\mu_{3,1}\circ Rb_{3,1*}a_{3,1}^*(\mu_{2,1}\boxtimes1)\circ C_L,\\
 M_R&=\mu_{2,2}\circ Rb_{2,2*}a_{2,2}^*(1\boxtimes\mu_{1,1})\circ C_R.
\end{split}
\]
These are maps of coefficient complexes, before cohomology.

\begin{theorem}[Associativity for block sizes $(2,1,1)$]
\label{ts22:thm:assoc211}
With the support comparison $\TS_{f,g}$ and the complex orientations
above, $M_L=M_{211}=M_R$ in the rational constructible derived
category on $\frM_4$. Consequently
\[
 m_{3,1}(m_{2,1}(x,y),z)=m_{2,2}(x,m_{1,1}(y,z))
\]
for $x\in H^*(\frM_2,\cP_2)$ and $y,z\in H^*(\frM_1,\cP_1)$.
Both sides commute with potential monodromy and with specialization
after pullback to an analytic family of commuting triples.
\end{theorem}

\begin{proof}
In a frame ordered by blocks $(2,1,1)$, the forbidden lower entries
are the block entries $(21),(31),(32)$, with sizes two, two, one,
for each of the three matrices. The two factorizations group the
same fifteen complex normal directions differently. Their Thom
classes multiply to the complete normal class. Every normal degree
is even, so interchanging these factors gives sign $+1$.
The flag variety is a Grassmannian bundle over the three-space
Grassmannian in the first factorization and a projective-line
bundle over the plane Grassmannian in the second. Iterated
integration in both cases evaluates on the same complex fundamental
cycle of the five-dimensional flag variety. Thus
\[
 \gamma_{3,1}\,Rq_{3,1*}\gamma_{\mathrm{ref},L}
       =\gamma_{211}
       =\gamma_{2,2}\,Rq_{2,2*}\gamma_{\mathrm{ref},R}.
\]

Smooth transverse base change identifies the refinement Gysin maps
with the indicated smaller incidences on the ambient spaces.
Proposition~\ref{ts22:prop:coherence} commutes their external coefficient
products, support comparisons, and proper direct images. In
particular the right side uses the natural square for
$1\boxtimes\gamma_{1,1}$: its source is supported where the second
potential is zero, but its target has the genuine second function
$f_2$. Naturality of the inclusion of supports is the required
comparison between these two coefficient operations.
Applying $\Phi_{f_4}$ and restricting by proper base change therefore
identifies the two displayed ambient composites with $M_L$ and
$M_R$, respectively, and their common value with $M_{211}$.
The identifications are constructed on the action diagrams, so the
equality descends to the quotient stacks.

The tensor associator preserves the order $(x,y,z)$ and introduces
no permutation sign. The displayed Gysin shifts are
even. This proves the stated identity on cohomology, also when $x$
has odd degree. Rotating the common half-plane commutes with the
support inclusions; proper exchange and Thom integration retain
this compatibility. Pulling the equality back to a family and
applying its central-restriction-to-nearby-cycles map proves the
specialization assertion. Vanishing cycles of the ambient potential
are pulled back here; they are not recomputed from its zero
restriction on the family.
\end{proof}

\section{The deciding coefficient groups and signs}
\label{ts22:sec:stalks}

The comparison retains the extra cubic groups at a rank-two scalar
triple. To calculate them, the nonzero fibre of $2\det M$ can be scaled to $\SL_3(\C)$. Polar decomposition
$M=UP$, with $U\in\SU_3$ and $P$ positive Hermitian of determinant
one, gives the retraction $UP\mapsto UP^t$, $1\ge t\ge0$.
It stays inside a bounded determinant fibre: if the eigenvalues of
$P$ are $e^{x_i}$ with $\sum x_i=0$, then
$\|UP^t\|^2=\sum e^{2tx_i}$ is convex in $t$ and has derivative
zero at zero, so it decreases toward $t=0$. For a sufficiently small
nonzero determinant value the whole compact $\SU_3$ lies in the
scaled ball. Thus the bounded local fibre retracts onto $\SU_3$.

The first-column map $\SU_3\to S^5$ has fibre $\SU_2=S^3$;
Gram--Schmidt extension of a column supplies local trivializations.
Cover $S^5$ by two enlarged hemispheres. The total spaces retract
onto $S^3$ and the overlap onto $S^4\times S^3$. The difference maps
on degrees zero and three are $(a,b)\mapsto a-b$, since the transition
preserves the fibre orientation. The cochain exact sequence of the
cover then gives groups in degrees $0,3,5,8$, each of rank one.
The degree-three fibre class and the pulled-back degree-five sphere
class have product evaluating nontrivially on the total orientation,
by iterated fibre integration. Hence
\[
 H^*(\SU_3;\Q)=\Lambda(\alpha_3,\alpha_5),
 \qquad \alpha_8:=\alpha_3\alpha_5.
\]
The normalized rank-two scalar coefficient has generators in degrees
$-4,-2,1$, corresponding to reduced Milnor degrees $3,5,8$.
At a nonzero rank-one traceless triple the quadratic fibre retracts
onto $S^3$, so there is only the degree-$-4$ group.

For two scalar rank-two blocks, Lemma~\ref{ts22:lem:wall} and the
relative-cochain product compute the complete stalk of the sum
coefficient on the product ambient space. Write
$[\alpha_p,\alpha_q]=\TS_{f_2,f_2}(\alpha_p\otimes\alpha_q)$.
This is a defined cochain product, not merely a matching of ranks.
Its groups are
\[
\begin{array}{c|rrrrrr}
 \text{normalized degree}&-8&-6&-4&-3&-1&2\\ \hline
 \text{rank}&1&2&1&2&2&1\\
 \text{reduced Milnor degree of }f_2+f_2&7&9&11&12&14&17.
\end{array}
\]
Indeed a pair $(p,q)$ contributes to reduced degree $p+q+1$ and
normalized degree $p+q-14$. The nine ordered pairs with
$p,q\in\{3,5,8\}$ give the displayed numbers.
If precisely one block is scalar and the other has nonzero rank-one
traceless part, the normalized degrees are $-8,-6,-3$, each of rank
one. If neither is scalar, only degree $-8$ remains; its Milnor
fibre is the sphere $S^7$ of an eight-variable nondegenerate quadratic.

The interchange of the two blocks acts on these products by
\[
 [\alpha_p,\alpha_q]\longmapsto
       (-1)^{(p-7)(q-7)}[\alpha_q,\alpha_p].
\]
In particular $[\alpha_8,\alpha_8]$ changes sign. Its nonzero value
in degree two is a decisive check that the comparison retains the
odd scalar coefficient. It cannot be replaced by an ungraded product.
The shifts eight and sixteen are even, so this is exactly the
Koszul sign of the normalized complexes.

Potential monodromy on the determinant fibre acts by multiplication
by $e^{2\pi i/3}$. Left multiplication by
\[
 \operatorname{diag}(e^{2\pi it/3},e^{2\pi it/3},e^{-4\pi it/3}),
 \qquad 0\le t\le1,
\]
is a norm-preserving homotopy on $\SL_3$ from the identity to that
transformation. It proves identity on these local cohomology groups.
Quadratic monodromy is the antipodal map of $S^3$, of degree $+1$.
The support comparison thus gives identity monodromy on every group
in the displayed table. This assertion is about these stalk groups;
it does not split the ambient complexes into their cohomology sheaves.

\section{Four distinct eigenvalues and the geometric boundary}
\label{ts22:sec:regular}

Suppose $A\in\Mat_4$ has four distinct eigenvalues. The commuting
matrices $B,C$ are diagonal in its eigenbasis. Off the diagonal,
\[
 f_4=\sum_{i<j}(a_i-a_j)
                 (b_{ij}c_{ji}-b_{ji}c_{ij})
\]
in a transverse slice; the eigenbasis directions and diagonal entries
are free. There are six pairs, each contributing four nondegenerate
normal variables. The twenty-four-variable quadratic gives
$\cP_4=\Q[8]$ on this regular locus. Each rank-two block contributes
$\Q[4]$, so the incoming comparison has the same shift $[8]$.
The real normal orientations can be fixed by their complex paired
$B,C$ coordinates. The product support comparison carries the product
normal generator to that of their sum, by the product of relative
normal balls.

There are six invariant planes, one for each two-element subset of
the four eigenlines. At any such plane the normal incidence is a
maximal isotropic subspace in the complementary hyperbolic directions.
Its Thom boundary has coefficient one: projection to the complementary
normal coordinates gives the angular sphere of integral one.
Interchanging the order of two eigenlines exchanges two hyperbolic
pairs and has sign $+1$. Hence $\mu_{2,2}$ is the sum of these six
sheet coefficients. This fixes the regular target generator using
complex Thom orientations and the support product simultaneously.

For flags of type $(2,1,1)$ there are twelve sheets. The left
parenthesization chooses a three-space in four ways and a plane
inside it in three ways. The right chooses a plane in six ways
and orders the remaining two lines in two ways. Both maps are the
same sum on twelve coordinates, giving $4\cdot3=6\cdot2=12$ on
the invariant constant class. This is a computation of the maps
already compared in Theorem~\ref{ts22:thm:assoc211}.

At a scalar four-dimensional module, the plane fibre is
$\operatorname{Gr}(2,4)$ and the incoming coefficient is a family
of two scalar rank-two coefficient complexes on its tautological
subspace and quotient. The nine-group calculation in
Section~\ref{ts22:sec:stalks} describes its individual ordinary atlas
stalks. It does not trivialize that family or remove its frame
actions. The map $\mu_{2,2}$ uses the whole coefficient family
before proper direct image. The full scalar stalk of $\cP_4$, its
image under this pushforward, and its extension structure require
additional computations; no rank for that target is asserted here.

For point modules, eigenvalues are their positions in $\C^3$ and
nilpotent parts describe extensions at coincident positions. Two
rank-two constituents can each have cubic critical fluctuations.
Their sum is therefore governed by the product of their critical
complexes, including the odd scalar class and its interchange sign.
The oriented flag map couples these coefficients to the four-dimensional
module. The theorem proves agreement of the two assemblies of
sizes $(2,1,1)$ in ordinary rational critical cohomology. Products
in further block sizes, higher coherences, a coproduct, Tate
normalizations, and a compact-threefold field-theoretic comparison
remain separate constructions.
